\section{Projects}
    {\textcolor{linkColor}{\textit{\textbf{Real-Time Watch Together Platform (\href{https://github.com/ahmniab/ModdieView.git}\textbf{ModdieView \faIcon{external-link-alt}}}}, \href{https://github.com/ahmniab/Moddie-View-Backend}{\textcolor{linkColor}{\textbf{ModdieView Backend \faIcon{external-link-alt}}}})}}
\vspace{0.10 cm}
\begin{onecolentry}
   \textit{A \textbf{real-time, synchronized media viewing application} to watch videos together with friends. }
    \begin{highlights}
        \item Implemented a a support for \textbf{Youtube}, \textbf{Vimoe} and \textbf{Direct Links}.
        \item Integrated the \textbf{YouTube Data API} for in-room video search.
        \item Built real-time synchronization for playback controls (play, pause, seek, and playback speed) using \textbf{Socket.IO}.
        \item Developed real-time chat with replies, reactions, typing indicators, message actions, and notifications.
        \item Implemented room management with invite links and multi-user interactions.
    \end{highlights}
\end{onecolentry}

\vspace{0.5 cm}
\newpage

    % \textbf{Planets Diseases} \hfill 
    \textit{\href{https://github.com/ahmniab/planets-diseases.git}{\textcolor{linkColor}{\textbf{Planets Diseases \faIcon{external-link-alt}}}}}
    

\vspace{0.10 cm}
\begin{onecolentry}

\textit{A full-stack web application built with \textbf{Next.js 16} and 
\textbf{TypeScript} that allows users to browse plant diseases and enables 
admins to manage structured documentation using \textbf{Firebase (Auth \& Firestore)} 
with optional \textbf{Redis} caching for performance optimization.}

\vspace{0.05 cm}

\begin{highlights}
    
    \item Developed a rich documentation system using \textbf{Editor.js} supporting structured content (headers, tables, images, lists).
    \item Built responsive UI with \textbf{Material-UI (MUI v7)} and managed client state using \textbf{TanStack Query v5}.
    \item Integrated \textbf{React Hook Form + Yup} for validation and \textbf{Axios} for API communication.
\end{highlights}

\end{onecolentry}
\vspace{0.5 cm}

% so what do u think?
% 
\textit{\href{https://github.com/ahmniab/SportsTournamentScheduling}{\textcolor{linkColor}{\textbf{Sports Tournament Scheduling \faIcon{external-link-alt}}}}}

\vspace{0.10 cm}
\begin{onecolentry}
    \textit{A distributed microservices system built with \textbf{.NET Core}, \textbf{Python}, and \textbf{React} to manage sports tournament resources and automate optimal timetable generation.}
    \begin{highlights}
        \item Developed a Clean \textbf{Architecture microservices} backend using ASP.NET Core, utilizing \textbf{gRPC} for synchronous inter-service communication and \textbf{RabbitMQ} for event-driven orchestration.
        \item Configured distributed state management and caching using Redis to \textbf{queue, track, and persist active generation job metadata} across C\# and Python runtimes.
        \item Integrated Keycloak for Identity and Access Management \textbf{(OIDC)} alongside \textbf{SpiceDB} to support Zanzibar-inspired fine-grained \textbf{relationship authorization}.
        \item Created a \textbf{React/TypeScript} frontend dashboard featuring live updates for active background scheduling pipeline jobs.
    \end{highlights}
\end{onecolentry}

\vspace{0.5 cm}

\textit{\href{https://github.com/ahmniab/ecg-detection-system}{\textcolor{linkColor}{\textbf{ECG Detection System \faIcon{external-link-alt}}}}}

\vspace{0.10 cm}
\begin{onecolentry}
    \textit{A \textbf{real-time}, \textbf{microservices-based electrocardiogram (ECG)} monitoring and anomaly detection system that ingests high-frequency cardiac signals, runs \textbf{AI-driven arrhythmia classifications}, and visualizes live telemetry.}
    \begin{highlights}
        \item Engineered a \textbf{high-throughput Go} ingestion service utilizing \textbf{gRPC streaming} and TimescaleDB bulk-copy protocols to performantly persist high-frequency voltage metrics.
        \item Architected a decoupled, \textbf{event-driven data pipeline using Apache Kafka} to stream heart rate patches to a \textbf{Python AI} classifier for real-time arrhythmia detection.
        \item Implemented a secure \textbf{Backend-for-Frontend (BFF)} pattern in ASP.NET Core with \textbf{OpenID Connect (OIDC)} and Keycloak to authenticate clients and manage session credentials.
        \item Developed an interactive React SPA featuring \textbf{Apache ECharts} to performantly plot thousands of \textbf{timeseries ECG samples} and overlay diagnosed cardiac events.
    \end{highlights}
\end{onecolentry}
\vspace{0.5 cm}

\textit{\href{https://github.com/ahmniab/Automated-E-Commerce-Deployment-Platform}{\textcolor{linkColor}{\textbf{Automated E-Commerce Deployment Platform \faIcon{external-link-alt}}}}}

\vspace{0.10 cm}
\begin{onecolentry}
    \begin{highlights}
        \item Refactored an existing monolithic codebase by decoupling it into an independent microservices architecture to improve maintainability and scalability.
        \item Containerized the separated microservices using Docker and orchestrated the multi-container environment utilizing Docker Compose for seamless inter-service communication.
        \item Designed and implemented an automated CI/CD pipeline using Jenkins to streamline integration and deployment workflows.
        \item \textbf{Tools:} Jenkins, Docker, Docker Compose, Microservices Architecture
    \end{highlights}
\end{onecolentry}
\vspace{0.5 cm}


%     \textit{\href{https://github.com/ahmniab/GraphicsSln}{\textcolor{linkColor}{\textbf{Graphics Algorithms Web \faIcon{external-link-alt}}}}}

% \vspace{0.10 cm}
% \begin{onecolentry}
%     \begin{highlights}
%         \item Interactive web application demonstrating computer graphics algorithms
%         \item Implemented DDA and Bresenham's line drawing algorithms
%         \item Tools: ASP.NET, JavaScript, HTML5 Canvas
%     \end{highlights}
% \end{onecolentry}
% \vspace{0.5 cm}

%     \textit{\href{https://github.com/ahmniab/Projectile-motion-simulator}{\textcolor{linkColor}{\textbf{Projectile Motion Simulator \faIcon{external-link-alt}}}}}

% \vspace{0.10 cm}
% \begin{onecolentry}
%     \begin{highlights}
%         \item Developed a C-based physics simulation application using the Raylib graphics library to visualize projectile motion trajectories with customizable parameters (launch angle, initial velocity, gravity, and time scaling).
%         \item Implemented HD video rendering functionality by integrating FFmpeg for frame capture and encoding, enabling users to export high-quality (1920×1080) simulation videos with real-time preview and keyboard-controlled interactions.
%         \item Designed a modular, cross-platform architecture with separate modules for physics engine, UI controls (textboxes, checkboxes), and object management; built for Linux with Windows support planned, demonstrating proficiency in low-level C programming, build systems, and external library integration.
%     \end{highlights}
% \end{onecolentry}
